\documentclass[10pt,a4paper]{amsart}
\usepackage[margin=19mm]{geometry}
\usepackage{amsmath,amssymb,mathtools}
\usepackage[scr=rsfs,cal=euler]{mathalfa}
\usepackage{xfrac}
\usepackage[unicode,hidelinks,bookmarks=false]{hyperref}
\newcommand{\ie}{i.e.\ }
\newcommand{\cf}{cf.\ }
\newcommand{\resp}{respectively\ }
\newcommand{\Subsection}{\subsection}
\providecommand{\nobreakdash}{\nobreak\hbox{-}}
\setlength{\emergencystretch}{2em}
\multlinegap0pt
\usepackage{bm}
\usepackage{bbm}
\usepackage{dsfont}
\usepackage{xspace}
\usepackage{cancel}
\usepackage{mathtools}
\usepackage{amsthm}
\makeatletter
\@ifundefined{theorem}{\newtheorem{theorem}{Theorem}}{}
\@ifundefined{lemma}{\newtheorem{lemma}[theorem]{Lemma}}{}
\makeatother
\newcommand\cG{\mathcal G}
\newcommand\cV{\mathcal V}
\providecommand{\N}{}
\title[Achievable Burning Densities of Growing Grids]{Achievable Burning Densities of Growing Grids}
\author{Jordan Barrett, Karen Gunderson, JD Nir, Pawel Pralat}
\date{Source: arXiv:2601.14151v1 (CC BY 4.0)}
\begin{document}
\maketitle

\noindent\emph{Source and licence.} Achievable Burning Densities of Growing Grids. Version arXiv:2601.14151v1 (CC BY 4.0), distributed under the Creative Commons Attribution 4.0 International licence, as recorded in the arXiv licence metadata for this exact version and shown on the abstract page https://arxiv.org/abs/2601.14151v1. Full rights notice: licenses/arxiv-2601.14151-CC-BY-4.0.txt.\par\smallskip

\noindent\textit{Editorial scope.} Counted unit: the Lemma whose source label is \texttt{lem:f,g same densities} in the article (source0.tex lines 412--418, source label \texttt{lem:f,g same densities}), delivered together with the article's own complete original proof of it (source0.tex lines 420--482, one \texttt{proof} environment of 63 lines). No mathematics is written, repaired or re-derived: statement and proof are the article's own text line for line. Required context retained uncounted: none. Every definition, display and label of those lines is retained unchanged. Omitted from this package and not redistributed: 1--411, 419--419, 483--1027. Nothing inside a retained excerpt is omitted or altered: every retained body is delivered exactly as the article's lines are written, so no omission receipt of any kind accompanies this package. Citations: the retained excerpts carry no citation command at all, so no bibliography entry of the article is transcribed and no reference is redistributed. Reference adaptation: none was needed and none was made; every reference inside the delivered unit resolves inside it, so no unresolved reference or citation is produced. Printed slips kept literal: no printed word, symbol, hypothesis or number of the counted statement or of its proof was altered. Layout and class: the article's own class is \texttt{article} with packages including fullpage, ulem, amsthm, amsmath, amssymb, appendix, comment, enumerate, graphicx, mathrsfs, mathtools, thm-restate, xcolor, cleveref; the wrapper is the lane's pinned \texttt{amsart} editorial wrapper at 10pt on A4, which restates the theorem-like environments the retained text needs and adds the article's own macro definitions that the retained text needs (\texttt{\textbackslash N}, \texttt{\textbackslash cG}, \texttt{\textbackslash cV}), with extra wrapper packages (bm, bbm, dsfont, xspace, cancel, mathtools). The delivered numbering is the unit's own. Assets: no retained span contains a figure environment, a graphics-inclusion command, a PDF-page inclusion, a table or a TikZ picture; no image file is copied into this package, the package declares no asset dependency and no image file is redistributed with it. Licence: version arXiv:2601.14151v1 of this article is distributed under the Creative Commons Attribution 4.0 International licence, as recorded in the arXiv licence metadata for this exact version and shown on the abstract page https://arxiv.org/abs/2601.14151v1. A full rights notice is delivered at licenses/arxiv-2601.14151-CC-BY-4.0.txt. Characters: every retained excerpt is pure ASCII, so the delivered engine, XeLaTeX with its default Latin Modern text font, reports no missing character.
\begin{lemma}\label{lem:f,g same densities}
Let $f,g$ be increasing functions such that at least one of $f,g$ has controlled growth and $\lim_{n \to \infty} f(n)/g(n) = 1$. Then 
\begin{enumerate}
\item[(a)] $f$ and $g$ both have controlled growth, and
\item[(b)] $P(\cG(f)) = P(\cG(g))$. 
\end{enumerate}
\end{lemma}

\begin{proof}
Beginning with (a), since $f$ and $g$ are increasing they both satisfy condition~(i) of the controlled growth requirement. Now suppose $f$ has controlled growth and let $\epsilon: \N \to \N$ satisfy $\epsilon
(n) = o(n)$. Then
\begin{align*}
g(n+\epsilon(n)) 
&= 
(1+o(1)) f(n+\epsilon(n))\\
&=
(1+o(1)) f(n)\\
&=
(1+o(1)) g(n) \,.
\end{align*}
Therefore, $g$ satisfies condition~(ii) of the controlled growth requirement. 

Now fix $c > 0$ and let $d > 0$ be such that for sufficiently large $n$, $f(n+cn) \geq (1+d)f(n)$. Then
\begin{align*}
g(n+cn) 
&= 
(1+o(1)) f(n+cn)\\
&\geq 
(1+o(1)) (1+d)f(n)\\
&=
(1+d+o(1))g(n) \,.
\end{align*}
Thus, for sufficiently large $n$, $g(n+cn) \geq (1+d/2) g(n)$, meaning $g$ satisfies condition~(iii) of the controlled growth requirement. Therefore, if $f$ has controlled growth then so does $g$ and, by a symmetrical argument, if $g$ has controlled growth then so does $f$. 

Continuing with (b), write $\cG(f) = (G^{(f)}_n,n \geq 1)$ and $\cG(g) = (G^{(g)}_n,n \geq 1)$. Let $\rho \in P(\cG(f))$ and let $\cV = (v_n,n\geq 1)$ be an activator sequence on $\cG(f)$ such that $\delta(\cG(f),\cV) = \rho$. Define the activator sequence $\cV' = (v_n', n \geq 1)$ on $\cG(g)$ inductively as follows. First, if $v_1 \in V(G^{(g)}_1)$ then set $v_1'=v_1$, and otherwise set $v_1' = v_0 := (0,0)$. Then, for $n \geq 1$ and $k \geq 0$, if $v_n' = v_{n-k}$, set $v_{n+1}' = v_{n-k+1}$ if $v_{n-k+1} \in V(G^{(g)}_{n+1})$, and otherwise set $v_{n+1}' = \emptyset$. In words, $\cV'$ burns the same vertices as $\cV$ in the same order, except that $\cV'$ ``waits'' until it is able to burn the correct vertices in $\cG(g)$. We claim that $\underline{\delta}\big( \cG(g), \cV' \big) = \overline{\delta}\big( \cG(g), \cV' \big) = \rho$. 

Firstly, note that $\cV'$ is a well defined activator sequence on $\cG(f)$ as well as on $\cG(g)$. Moreover, after any turn $t$, the number of vertices burned by $\cV'$ in $\cG(f)$ is at most the number of vertices burned by $\cV$. Thus, $\overline{\delta}\big( \cG(f), \cV' \big) \leq \delta\big( \cG(f), \cV \big) = \rho$. Next, since $|V(G^{(g)}_n)| = (1+o(1))|V(G^{(f)}_n)|$, and since the sets of burning vertices in $(\cG(f),\cV')$ and in $(\cG(g),\cV')$ are identical at all times, we get that 
\[
\overline{\delta}\big( \cG(g), \cV' \big) = \overline{\delta}\big( \cG(f), \cV' \big) \leq \delta\big( \cG(f), \cV \big) = \rho \,.
\]

We are left to show that $\underline{\delta}\big( \cG(g), \cV' \big) \geq \rho$. For each $n \geq 1$, define $\epsilon(n)$ to be the smallest non-negative integer that satisfies $g(n+\epsilon(n)) \geq f(n)$. We claim that $\epsilon(n) = o(n)$. Suppose, to the contrary, that there exists a constant $c > 0$ and an increasing sequence $(n_i,i\geq 1)$ such that $\epsilon(n_i) \geq cn_i$ for all $i \geq 1$. Then, for each $i \geq 1$, by the minimality of $\epsilon$ we have $g(n_i+\epsilon(n_i)-1) < f(n_i)$. By condition~(ii) of the controlled growth requirement, this inequality implies that $g(n_i+\epsilon(n_i)) = (1+o(1)) f(n_i)$, implying further that 
\[
g(n_i + cn_i) \leq g(n_i+\epsilon(n_i)) = (1+o(1))f(n_i) = (1+o(1))g(n_i) \,.
\] 
However, by condition~(iii) of the controlled growth requirement, we can find a constant $d > 0$ such that for sufficiently large $n$, $g(n + cn) \geq (1+d)g(n)$. Therefore, the assumption that $\epsilon(n) \neq o(n)$ leads to a contradiction, and thus $\epsilon(n) = o(n)$.

Lastly, by the definition of $\epsilon$, the number of unique activator vertices in $(v_i,i \in [n]) \subseteq \cV$ that are not in $(v_i',i \in [n])$ is at most
\[
\epsilon_{\max}(i) := \max_{i \in [n]} \epsilon(i) \,,
\]
and thus the number of burning vertices in $(\cG(g),\cV')$ by the end of turn $t$ is at least the number of burning vertices in $(\cG(f),\cV)$ by the end of turn $t-\epsilon_{\max}(t) = (1-o(1))t$. Letting $B_t$ and $B_t'$ be the respective number of burning vertices in $(\cG(f),\cV)$ and $(\cG(g),\cV')$ by the end of turn $t$, we have that 
\begin{align*}
\frac{|B_t'|}{|V(G^{(g)}_t)|}
&\geq
\frac{|B_{t-\epsilon_{\max}(t)}|}{|V(G^{(g)}_t)|}\\
&= 
(1+o(1))\frac{|B_{t-\epsilon_{\max}(t)}|}{|V(G^{(f)}_t)|}\\
&= 
(1+o(1))\frac{|B_{t-\epsilon_{\max}(t)}|}{|V(G^{(f)}_{t-\epsilon_{\max}(t)})|} \,,
\end{align*}
the last equality holding since $f$ satisfies condition~(i) of the controlled growth requirement. Therefore,
\[
\liminf_{t \to \infty} \frac{|B_t'|}{|V(G^{(g)}_t)|}
\geq
\liminf_{t \to \infty} \frac{|B_t|}{|V(G^{(f)}_t)|} \,,
\]
implying $\underline{\delta}\big( \cG(g), \cV' \big) \geq \delta\big( \cG(f), \cV \big) = \rho$.

Therefore, $\delta \big( \cG(g), \cV' \big)$ is well defined and equals $\rho$. By symmetry, we have shown that $P(\cG(f)) = P(\cG(g))$, and this concludes the proof.  
\end{proof}

\end{document}
